\chapter{Conclusiones}
\label{cap:conclusiones}

Este cap\'itulo interpreta los resultados del \cref{cap:resultados} a la luz de los
objetivos y las hip\'otesis enunciados en los cap\'itulos anteriores. En primer lugar,
formula una conclusi\'on por cada objetivo espec\'ifico y precisa qu\'e permiten afirmar
los datos. A continuaci\'on, expone las l\'ineas de continuaci\'on abiertas por el
estudio.

Conviene advertir de entrada que el estudio es parcial. Las variantes V3 y V4 siguen en
ejecuci\'on, de modo que las conclusiones sobre el rendimiento se apoyan en las tres
variantes convolucionales completadas.

\section{Conclusiones del trabajo}
\label{sec:conclusiones}

\divisionmenor{Sobre la implementaci\'on de las variantes}

El primer objetivo espec\'ifico se cumple en su totalidad en cuanto a la implementaci\'on.
Las cinco variantes del codificador visual funcionan dentro de la misma
\textit{Diffusion Policy}, entregan un vector de condicionamiento de id\'entica
dimensi\'on y comparten red generativa, hiperpar\'ametros y protocolo de evaluaci\'on. La
comparaci\'on entre ellas queda as\'i acotada al bloque visual.

En cambio, la ejecuci\'on solo se ha completado para tres de las cinco variantes. Por
tanto, el trabajo responde ya a la segunda pregunta planteada en el \cref{sec:problema},
la relativa a la estrategia de entrenamiento, mientras que la primera, la relativa a la
arquitectura del codificador, queda abierta hasta disponer de V3 y V4.

\divisionmenor{Sobre el rendimiento de las variantes}

La conclusi\'on principal del trabajo contradice la hip\'otesis de partida: el bloque
formado por el codificador preentrenado sobre ImageNet y su cadena de preprocesado no
mejora el rendimiento en Push-T, sino que lo degrada de forma sustancial. La l\'inea base
entrenada desde cero alcanza una puntuaci\'on media de evaluaci\'on de \num{0.864}, frente
a \num{0.668} de la variante congelada y \num{0.648} de la ajustada. Las ca\'idas equivalen
al \SI{23}{\percent} y al \SI{25}{\percent} del valor de referencia. Ambas diferencias
superan la dispersi\'on entre episodios: sus intervalos de confianza al \SI{95}{\percent},
{[\num{0.091};\,\num{0.303}]} y {[\num{0.102};\,\num{0.332}]}, excluyen el cero, con
$p = \num{0.0004}$ y $p = \num{0.0001}$ en la prueba de los rangos con signo.

Asimismo, la segunda hip\'otesis tampoco se sostiene. El ajuste fino del codificador
preentrenado no supera a su congelaci\'on, puesto que V2 se queda \num{0.020} por debajo de
V1. Esa diferencia no alcanza significaci\'on ($p = \num{0.82}$) y su intervalo contiene el
cero, de modo que la conclusi\'on correcta no es que la congelaci\'on sea mejor, sino que
actualizar los pesos preentrenados no aporta ninguna ventaja medible. La incertidumbre
entre semillas de entrenamiento, no cuantificada, se suma a esa indeterminaci\'on.

El ajuste fino s\'i altera, en cambio, la din\'amica del entrenamiento, y siempre a peor.
V2 permanece estancada alrededor de \num{0.11} hasta la \'epoca \num{50}, cuando V0 y V1
ya rebasan \num{0.53}. Adem\'as, sobreajusta antes y con mayor intensidad: en la \'epoca
\num{250} su p\'erdida de validaci\'on alcanza \num{0.380}, frente a \num{0.120} de V1 y
\num{0.077} de V0. En conjunto, la actualizaci\'on de los pesos degrada la
representaci\'on de partida sin construir otra m\'as \'util con las demostraciones
disponibles.

Una explicaci\'on plausible atiende a la distancia entre dominios. Las observaciones de
Push-T son im\'agenes sint\'eticas de fondo uniforme, sin textura y con pocas formas,
mientras que las caracter\'isticas aprendidas sobre ImageNet codifican textura y categor\'ia
de objeto \cite{deng2009imagenet}. Esa informaci\'on resulta poco pertinente en una tarea
definida por la posici\'on relativa de dos figuras geom\'etricas. La observaci\'on es
coherente con dos antecedentes: R3M, preentrenado sobre v\'ideo, tampoco supera al
codificador entrenado de extremo a extremo en Push-T \cite{nair2022r3m},
\cite{chi2025diffusion}, y el estudio con DINOv3 concluye que la ventaja del
preentrenamiento depende de la tarea y de la estrategia de adaptaci\'on
\cite{egbe2025dinov3dp}.

La atribuci\'on causal exige, en cambio, una precisi\'on que recorre todo este apartado.
V0 difiere de V1 y V2 en cuatro factores a la vez, ya que el preprocesado viaja con los
pesos: la inicializaci\'on, la resoluci\'on de entrada (\num{96} frente a \num{224}
p\'ixeles), el recorte aleatorio, que solo aplica la l\'inea base, y la agregaci\'on
espacial a la salida del codificador. Sobre este \'ultimo conviene detenerse: V0 conserva
la localizaci\'on de las activaciones mediante un \textit{spatial softmax}
\cite{levine2016visuomotor}, mientras que las variantes preentrenadas entregan un
descriptor global. En una tarea definida por posiciones relativas, esa p\'erdida de
informaci\'on geom\'etrica puede resultar tan explicativa como la propia inicializaci\'on.

Las conclusiones anteriores se refieren, por tanto, al codificador junto con su cadena de
preprocesado, entendidos como un bloque, y no a la inicializaci\'on aislada. Separar los
cuatro efectos requiere entrenamientos adicionales, descritos en el
\cref{sec:trabajo-futuro}.

\divisionmenor{Sobre el coste computacional}

La l\'inea base no solo obtiene la mejor puntuaci\'on, sino que adem\'as es la m\'as
barata. Cada \'epoca de V0 ocupa \SI{1.5}{\minute}, frente a \SI{5.3}{\minute} de V1 y
\SI{14.7}{\minute} de V2, esto es, factores de \num{3.5} y \num{9.8}. En tiempo total, V0
consume \SI{17.4}{\hour} de GPU frente a \SI{96.5}{\hour} de V1 y \SI{84.9}{\hour} de V2,
y esta \'ultima cifra resulta menor solo porque su ejecuci\'on se detuvo antes. Las
variantes preentrenadas son, por tanto, simult\'aneamente peores y m\'as costosas,
combinaci\'on que no deja margen para un compromiso entre rendimiento y coste.

Ese sobrecoste no procede del n\'umero de par\'ametros entrenables. V1 congela el
codificador y entrena \num{277} millones de par\'ametros frente a los \num{288} millones de
V0, una reducci\'on del \SI{4}{\percent}, y aun as\'i cada \'epoca le cuesta tres veces y
media m\'as. El factor
determinante es la resoluci\'on de entrada, dado que una imagen de \num{224} p\'ixeles
contiene \num{5.4} veces m\'as p\'ixeles que una de \num{96}. De ah\'i se sigue una
conclusi\'on de alcance pr\'actico: congelar un codificador no abarata el entrenamiento
cuando obliga a ampliar la imagen de entrada.

\divisionmenor{Respuesta al objetivo general}

En conjunto, y para la configuraci\'on estudiada, la respuesta al objetivo general resulta
inequ\'ivoca en su parte convolucional. Con \num{206} demostraciones de Push-T, una
ResNet-18 entrenada de extremo a extremo a resoluci\'on nativa domina en rendimiento y en
coste a las dos alternativas preentrenadas sobre ImageNet. Ning\'un dato del estudio
respalda sustituirla por ellas.

Queda por determinar si esa conclusi\'on se extiende a los codificadores preentrenados con
objetivos distintos de la clasificaci\'on supervisada. Las representaciones
autosupervisadas de DINOv2 \cite{oquab2024dinov2} y las multimodales de CLIP
\cite{radford2021clip} podr\'ian conservar informaci\'on espacial m\'as pertinente para la
tarea. Los entrenamientos de V3 y V4 se encuentran en curso y decidir\'an ese punto.
[PENDIENTE: resultados de V3 y V4]

\section{Trabajo futuro}
\label{sec:trabajo-futuro}

Del estudio se derivan seis l\'ineas de continuaci\'on, ordenadas por su capacidad para
afianzar o corregir las conclusiones anteriores.

En primer lugar, completar las variantes V3 y V4. Ambas discriminan entre dos lecturas de
la conclusi\'on principal: que el problema resida en la supervisi\'on de ImageNet o que
afecte a todo codificador preentrenado fuera del dominio rob\'otico.

En segundo lugar, aislar los cuatro factores confundidos entre la l\'inea base y las
variantes preentrenadas. Un entrenamiento adicional de ResNet-18 desde cero a \num{224}
p\'ixeles y sin recorte aleatorio separa el efecto de la inicializaci\'on del efecto del
preprocesado. Su coste por \'epoca coincide con el de V2, del orden de
\SI{15}{\minute}, puesto que el codificador se actualiza en ambos casos y la resoluci\'on
de entrada es la misma; bajo el criterio de parada, el total se sit\'ua en torno a las
\SI{85}{\hour}. Una segunda ejecuci\'on desde cero a \num{96} p\'ixeles y sin recorte
aleatorio aislar\'ia el efecto del aumento de datos por unas \SI{17}{\hour}, coste
equivalente al de V0.

En tercer lugar, repetir las variantes seleccionadas con las semillas \num{43} y \num{44}.
Los contrastes del \cref{cap:resultados} cubren la dispersi\'on entre condiciones iniciales,
pero no la variabilidad entre ejecuciones de entrenamiento; solo la r\'eplica permite
estimarla, limitaci\'on declarada desde el \cref{sec:alcance} y cuantificada en el
\cref{sec:coste}.

En cuarto lugar, evaluar los puntos de control seleccionados sobre un bloque de semillas de
entorno disjunto del empleado para seleccionarlos. Ese paso convertir\'ia la medida actual,
que act\'ua como puntuaci\'on de selecci\'on, en una evaluaci\'on final independiente, y su
coste se reduce a tres ejecuciones de inferencia.

En quinto lugar, examinar la agregaci\'on espacial a la salida de los codificadores
preentrenados. La l\'inea base conserva la localizaci\'on de las activaciones mediante un
\textit{spatial softmax} \cite{levine2016visuomotor}, mientras que las variantes
preentrenadas entregan un descriptor global. Cabe evaluar si una agregaci\'on que preserve
la geometr\'ia recupera parte de la diferencia observada.

Por \'ultimo, extender el protocolo a otras tareas y a un robot f\'isico. Push-T aporta un
banco de pruebas controlado, pero sus im\'agenes sint\'eticas son precisamente el escenario
menos favorable para un codificador preentrenado sobre fotograf\'ias naturales. La
conclusi\'on podr\'ia invertirse en observaciones reales, extremo que este trabajo no puede
resolver.
